\documentclass[10pt,a4paper]{amsart}
\usepackage[margin=19mm]{geometry}
\usepackage{amsmath,amssymb,mathtools}
\usepackage[scr=rsfs,cal=euler]{mathalfa}
\usepackage{xfrac}
\usepackage[unicode,hidelinks,bookmarks=false]{hyperref}
\newcommand{\ie}{i.e.\ }
\newcommand{\cf}{cf.\ }
\newcommand{\resp}{respectively\ }
\newcommand{\Subsection}{\subsection}
\providecommand{\nobreakdash}{\nobreak\hbox{-}}
\setlength{\emergencystretch}{2em}
\multlinegap0pt
\usepackage{amssymb}
\usepackage{mathtools}
\usepackage{float}
\usepackage[shortlabels]{enumitem}
\newtheorem{lemma}{Lemma}
\DeclareMathOperator{\GL}{GL}
\DeclareMathOperator{\Sym}{\mathsf{S}^2}
\DeclareMathOperator{\diag}{diag}
\newcommand{\tp}{{\scriptscriptstyle\mathsf{T}}}
\DeclareMathOperator{\tr}{tr}
\title{Pierce--Birkhoff conjecture is false}
\author{Zehua Lai, Lek-Heng Lim, Junyu Ren}
\date{Source: arXiv:2609.10420v1 (CC BY 4.0)}
\begin{document}
\maketitle

\noindent\emph{Source and licence.} Pierce--Birkhoff conjecture is false. Version arXiv:2609.10420v1 (CC BY 4.0), distributed by arXiv under the Creative Commons Attribution 4.0 International licence, http://creativecommons.org/licenses/by/4.0/, as recorded in the arXiv licence metadata for this exact version and shown on the abstract page https://arxiv.org/abs/2609.10420v1. Full licence notice: \texttt{licenses/arxiv-2609.10420-CC-BY-4.0.txt}.\par\smallskip

\noindent\textit{Editorial scope.} Counted unit: the article's lemma lem:quadratic-orbit (source0.tex lines 209--215). The article's complete original proof of it is delivered with it (source0.tex lines 217--278, one \texttt{proof} environment of 62 lines). No mathematics is written, repaired or re-derived: the statement and the proof are the article's own lines, in the article's own notation, and no retained line is substituted or re-flowed.  No further source-local context is needed: every cross-reference of every retained line resolves inside the two delivered spans.  Nothing was omitted from any retained span, so the package carries no omission record.  No retained line contains a citation, so the wrapper carries no bibliography and no bibliography entry.  No retained line contains a figure environment, an image-inclusion command or drawing code, so no image is delivered and the package declares no asset dependency.  Editorial formatting only, in addition: an AMS \texttt{amsart} preamble that restates the article's own declarations and macros used by the retained lines, namely the environments \texttt{lemma} on separate counters, together with the standard packages those lines need.  The retained environments print in this document as Lemma 1 in order of appearance, whereas the article prints the same items with the numbering of its own full document.  The complete native source of the article as distributed in the arXiv e-print for this exact version is bundled unchanged as \texttt{sources/209910/source0.tex}.
\begin{lemma}[Quadratic orbit lemma]\label{lem:quadratic-orbit}
A real polynomial $p(X,Y)$ of total degree at most two vanishes on $\mathcal O$ if and only if
\[
p(X,Y)=L(XY)
\]
for some real linear functional $L$ on $\mathbb{R}^{5 \times 5}$.
\end{lemma}

\begin{proof}
The reverse implication is immediate from $XY=0$ on $\mathcal O$.

Now we prove that if $p$ vanishes on $\mathcal O$ and has degree at most 2, then $p$ has the desired form. First note that replacing $U$ by
\[
U\diag(a,a,b^{-1},b^{-1},1)
\]
replaces $(X,Y)$ by $(a^2X,b^2Y)$. Hence if $p$ vanishes on $\mathcal O$, then all bihomogeneous components of $p$ vanish on $\mathcal O$.

The projection of $\mathcal O$ to the first coordinate $X$ consists of positive semidefinite symmetric matrices of rank two. Its Euclidean closure contains every rank-one matrix $uu^\tp $. A linear form vanishing on all $uu^\tp $ is zero since matrices of the form $uu^\tp $ span the space $\Sym(\mathbb{R}^5)$.

If a homogeneous quadratic form $Q$ vanishes on the rank-two positive semidefinite locus, then it also vanishes on $uu^\tp $ and on $uu^\tp +vv^\tp $. Polarization gives
\[
\widetilde Q(uu^\tp ,vv^\tp )=0
\]
for all $u,v$. Again, since rank-one symmetric matrices span $\Sym(\mathbb{R}^5)$, $Q=0$. The same argument applies to pure $Y$ components. The constant component is also zero. Thus only a bilinear form $B(X,Y)$ may remain.

We next show that the Euclidean closure of $\mathcal O$ contains every pair
\[
(uu^\tp ,ss^\tp )\qquad\text{with }u^\tp s=0.
\]
Take
\[
U_\delta=\diag(1,\delta,1,\delta^{-1},1).
\]
Then
\[
U_\delta DU_\delta^\tp =E_{11}+\delta^2E_{22},
\qquad
U_\delta^{-\tp} EU_\delta^{-1}=E_{33}+\delta^2E_{44},
\]
which tends to $(E_{11},E_{33})$ as $\delta \to 0$. If $u^\tp s=0$, we can choose $W\in\GL_5(\mathbb R)$ with $We_1=u$ and $W^{-\tp} e_3=s$; applying the fixed change of basis $W$ gives the desired pair in the closure.

Hence
\[
B(uu^\tp ,ss^\tp )=0\qquad\text{whenever }u^\tp s=0.
\]
Set
\[
F(u,s)=B(uu^\tp ,ss^\tp ).
\]
This polynomial has bidegree $(2,2)$ and vanishes on the hypersurface $u^\tp s=0$. The polynomial $u^\tp s$ is irreducible, and its real zero set is Zariski dense in that hypersurface. Therefore
\[
F(u,s)=(u^\tp s)G(u,s)
\]
with $G$ of bidegree $(1,1)$. Thus $G(u,s)=u^\tp As$ for some matrix $A$, so
\[
F(u,s)=(u^\tp s)(u^\tp As).
\]
Define $L(M)=\tr(A^\tp M)$. Since
\[
(uu^\tp )(ss^\tp )=(u^\tp s)us^\tp ,
\]
we have
\[
B(uu^\tp ,ss^\tp )=L(uu^\tp ss^\tp ).
\]
Rank-one symmetric matrices span $\Sym(\mathbb{R}^5)$, so bilinearity extends the identity to all symmetric $X,Y$:
\[
B(X,Y)=L(XY). \qedhere
\]
\end{proof}

\end{document}
